\documentclass{article}

\usepackage{Paper}
\pdftitle{PAIND-description}

\title{\textbf{PAIND description}}
\author{Matteo Frongillo}
\date{}

\begin{document}
\maketitle

\begin{itemize}
    \item \textbf{Title}: Cold-water circulation with a heat pump
    \item \textbf{Supervisor}: Ludger Fischer
    \item \textbf{Industrial partner / CC}: R. Nussbaum AG / Patrik Zeiter
    \item \textbf{Description}:\\
        The project consists of a domestic water system with separate domestic cold water (DCW) and domestic hot water (DHW) distribution. The starting point of the project is the hygienic risk on the DCW side: during periods without flow, the DCW temperature can exceed the critical threshold of 20°C, creating conditions that may promote Legionella growth.

        To address this, a closed DCW loop is introduced before the domestic distribution. The actively cooled DCW loop directly contributes to hygiene by keeping the cold-water temperature below the critical range. The heat extracted from the DCW loop can be reused for DHW preparation through a heat pump.

        The DHW side includes a circulation loop to maintain hygiene and allow thermal disinfection against Legionella. The use of recovered heat for DHW production is a useful secondary benefit, but the main focus of the project is to assess whether active cooling of the DCW loop can improve hygienic safety while technically supporting heat pump operation.
\end{itemize}

\end{document}
